\subsection{闭子群空间拓扑的另一种解释}

设 F 为 G 的闭子集所成的集合；按如下方式在 F 上定义一个分离一致结构：对 G 的每个紧子集 K 以及 e 在 G 中的每个邻域 V，设 \(\mathrm{P}(\mathrm{K},\mathrm{V})\) 为 F 中同时满足下列两个条件的元素对 \((\mathrm{X},\mathrm{Y})\) 所成的集合：

\[
\mathrm{X} \cap \mathrm{K} \subset \mathrm{VY} \quad \text { and } \quad \mathrm{Y} \cap \mathrm{K} \subset \mathrm{VX}.\tag{9}
\]

我们来证明，诸 \(\mathrm{P}(\mathbf{K},\mathbf{V})\) 是一个分离一致结构 \(\mathcal{U}\) （在 \(\mathfrak{F}\) 上的）的邻偶基本系。GT, II, §1, No.~1 的公理 \((\mathrm{U}_{\mathrm{I}}^{\prime})\) 与 \((\mathrm{U}_{\mathrm{II}}^{\prime})\) 显然满足；此外，关系式 \(\mathbf{K} \subset \mathbf{K}'\) 与 \(\mathbf{V}' \subset \mathbf{V}\) 蕴含 \(\mathrm{P}(\mathbf{K}',\mathbf{V}') \subset \mathrm{P}(\mathbf{K},\mathbf{V})\) ；因此为验证 \((\mathrm{U}_{\mathrm{III}}^{\prime})\) ，可以只限于 \(\mathbf{V}\) 是 \(e\) 的对称紧邻域的情形，此时 VK 是紧的。设 \((\mathbf{X},\mathbf{Y}) \in \mathrm{P}(\mathrm{VK},\mathbf{V})\) 且 \((\mathbf{Y},\mathbf{Z}) \in \mathrm{P}(\mathrm{VK},\mathbf{V})\) ；那么 \(\mathbf{X} \cap \mathbf{K} \subset \mathbf{X} \cap \mathrm{VK} \subset \mathrm{VY}\) ，并且若 \(y \in \mathbf{Y}\) 使得 \(vy \in \mathbf{K}\) （对某个 \(v \in \mathbf{V}\) ），则必然 \(y \in \mathrm{VK}\) ，因此

\[
\mathrm{X} \cap \mathrm{K} \subset \mathrm{V} (\mathrm{Y} \cap \mathrm{VK});
\]

另一方面，\(Y \cap VK \subset VZ\) ，由此 \(X \cap K \subset V^{2}Z\) ，并且同理可证 \(Z \cap K \subset V^{2}X\) ，这就证明了 \((U_{III}^{\prime})\) 。最后，若 X、Y 是 F 的两个不同元素，则例如存在一点 \(a \in X\) 使得 \(a \notin Y\) ，从而存在 e 的一个对称紧邻域 V 使得 \(Va \cap Y = \varnothing\) ，即 \(a \notin VY\) ；更有 \((X,Y) \notin P(Va,V)\) ，这就完成了断言的证明。

这一点既已确立，我们在闭子群所成的集合 \(\Sigma\) （\(G\) 的）上考虑拓扑 \(\mathcal{T}\) ，它由刚刚定义的一致空间 \(\mathfrak{F}\) 的拓扑所诱导。我们将看到这一拓扑与 No.~3 中定义的拓扑相同。只需证明映射 \(\alpha \mapsto H_{\alpha}\) （\(\Gamma\) 到 \(\Sigma\) 中的映射）在 \(\Sigma\) 赋有拓扑 \(\mathcal{T}\) 时是连续的；因为这时该映射在 \(\Gamma_{\varphi}\) 上的限制（记号同 No.~1 命题 2）也是如此，而它是双射的；但由于 \(\Gamma_{\varphi}\) 紧且拓扑 \(\mathcal{T}\) 是分离的，映射 \(\alpha \mapsto H_{\alpha}\) （\(\Gamma_{\varphi}\) 到 \(\Sigma\) 中的）于是将是一个同胚。

于是设 \(\alpha_0\) 为 \(\Gamma\) 的一点，\(\Phi\) 为 \(\Gamma\) 上收敛于 \(\alpha_0\) 的一个滤子；我们要证明，关于 \(\Phi\) ，\(\mathrm{H}_{\alpha}\) 趋向于 \(\mathrm{H}_{\alpha_0}\) （对拓扑 \(\mathcal{T}\) 而言）。设 \(\mathbf{K}\) 为 \(\mathbf{G}\) 的紧子集，\(\mathbf{V}\) 为 \(e\) 在 \(\mathbf{G}\) 中的一个对称紧邻域；对每一点 \(x \in \mathrm{H}_{\alpha_0} \cap \mathrm{K}\) ，存在一个集合 \(\mathrm{M}(x) \in \Phi\) ，使得对每个 \(\alpha \in \mathrm{M}(x)\) 都有 \(\mathrm{V}x \cap \mathrm{H}_{\alpha} \neq \emptyset\) (No.~1, Lemma 2)，从而 \(\mathrm{V}x \subset \mathrm{V}^2\mathrm{H}_{\alpha}\) ；把 \(\mathrm{H}_{\alpha_0} \cap \mathrm{K}\) 用有限个集合 \(\mathrm{V}x_i\) 覆盖，就可以看出，若 \(\mathrm{M} = \bigcap \mathrm{M}(x_i)\) ，则 \(\mathrm{H}_{\alpha_0} \cap \mathrm{K} \subset \mathrm{V}^2\mathrm{H}_{\alpha}\) 对每个 \(\alpha \in \mathrm{M}\) 都成立。

反之，假设存在 e 在 G 中的一个开邻域 U，使得对每个集合 \(L \in \Phi\) ，至少有一个 \(\alpha \in L\) 满足 \(H_{\alpha} \cap K \not\subset UH_{\alpha_{0}}\) ；若 \(\omega(L)\) 是具有这一性质的 \(\alpha \in L\) 所成的集合，则诸 \(\omega(L)\) 构成滤子 \(\Phi'\) （\(\Gamma\) 上比 \(\Phi\) 更细的）的滤基，并且对每个 \(\alpha\) （它属于诸 \(\omega(L)\) （\(L \in \Phi\)）的并 E ），存在一个 \(t_{\alpha} \in H_{\alpha} \cap K\) ，它不属于 \(UH_{\alpha_0}\) ；对 \(\alpha \notin E\) ，取 \(t_{\alpha}\) 为 \(H_{\alpha}\) 的任一点。由于 \(K \cap \mathbb{C}(UH_{\alpha_0})\) 是紧的，将存在一个聚点 \(s\) （\(\alpha \mapsto t_{\alpha}\) 关于 \(\Phi'\) 的），它属于 \(K \cap \mathbb{C}(UH_{\alpha_0})\) ；但由于 \(\Phi'\) 收敛于 \(\alpha_0\) （在 \(\Gamma\) 中），这与 No.~1 的引理 3 矛盾。

